\documentclass[journal]{IEEEtran}
\usepackage{amsmath,amssymb}
\usepackage{booktabs}
\usepackage{graphicx}
\usepackage{url}
\usepackage{xcolor}
\usepackage{hyperref}
\usepackage{enumitem}
\hypersetup{hidelinks}
\graphicspath{{figures/}}
\newcommand{\placeholder}[1]{\textcolor{blue}{\textit{[#1]}}}

\title{AI-Powered Real-Time Detection and Prevention of Voice Cloning Impersonation Attacks in VoIP Calls}

\author{Muthu Amudhan S, Amanjeeta Rath, Khushi Kumari, and Debosmita Mukhopadhyay%
\thanks{The authors are with \placeholder{Department and Institution, City, Country}. E-mail: \texttt{<muthu-email-placeholder>}, \texttt{<amanjeeta-email-placeholder>}, \texttt{<khushi-email-placeholder>}, and \texttt{<debosmita-email-placeholder>}.}}

\begin{document}
\maketitle

\begin{abstract}
Voice-cloning systems increasingly enable synthetic speech to impersonate a trusted caller in telephone and Voice-over-IP (VoIP) interactions. This paper presents a privacy-preserving architecture for real-time detection of synthetic speech using a classical digital signal processing (DSP) fast path and an extensible deep-learning path. The system validates and processes 16-kHz mono speech in three-second sliding windows, applies voice activity detection, extracts Fourier, spectral, phase, pitch, and speech-activity features, and combines calibrated evidence using a dependence-aware sequential Bayesian risk engine. The deep-learning interface supports a convolutional baseline and an XLS-R encoder with learned sensitive-layer selection and an AASIST-style graph-attention backend. We describe calibration, grouped zero-day evaluation, privacy-preserving telemetry, asynchronous inference batching, and GPU-oriented deployment. Quantitative results are intentionally reserved for experiments on the authors' labeled dataset and are not fabricated in this draft.
\end{abstract}

\begin{IEEEkeywords}
voice deepfake detection, synthetic speech, voice cloning, VoIP security, digital signal processing, XLS-R, calibration, sequential Bayesian inference.
\end{IEEEkeywords}

\section{Introduction}
Voice cloning can reproduce speaker identity cues from short recordings and can therefore be used to impersonate a caller during authentication, financial, or social-engineering interactions. Telephone channels complicate detection because codecs, bandwidth limitation, packet loss, resampling, and noise can create artifacts that resemble synthetic-speech artifacts. A practical detector must operate under latency constraints without treating ordinary channel artifacts as definitive evidence of synthesis.

This work develops a modular detector with four objectives: 1) provide a reproducible DSP baseline; 2) define a trainable deep-learning detector interface; 3) accumulate evidence across streaming windows without treating overlapping observations as independent; and 4) expose a deployable security service for authorized VoIP applications. The system processes audio in memory and emits metadata-only security events.

The contributions are:
\begin{itemize}[leftmargin=*]
\item a configurable three-second streaming pipeline with VAD gating and bounded ephemeral buffering;
\item a Fourier- and DSP-based feature engine that handles narrowband telephony limitations;
\item CNN and XLS-R/AASIST-compatible detector interfaces with standardized logits and embeddings;
\item calibrated score fusion and discounted sequential Bayesian risk accumulation;
\item grouped evaluation protocols for speaker-, source-, generator-, codec-, accent-, and noise-disjoint experiments; and
\item a FastAPI, Docker, MLflow, ZenML, and AWS-oriented operational foundation without raw call-audio persistence.
\end{itemize}

\section{Threat Model and Scope}
We consider an attacker who supplies or generates synthetic speech intended to impersonate a human speaker over an authorized telephone or VoIP channel. The detector observes audio and optional non-sensitive session identifiers. It does not assume caller identity, transaction context, or physiological ground truth. A low-risk call continues, a medium-risk call receives a verification recommendation, and a high-risk call produces an alert recommending step-up authentication. Automatic call termination or transaction freezing is outside the scope of this version.

The input is a mono floating-point waveform sampled at 16 kHz after channel-specific preprocessing. The reference evaluation unit is a three-second speech window with a configurable one-second hop. Labels follow the dataset convention $y=0$ for a human recording and $y=1$ for an AI-generated recording.

\section{System Architecture}
Fig.~\ref{fig:architecture} summarizes the processing flow. Audio is validated, resampled when required, held in a bounded RAM buffer, gated by VAD, and divided into overlapping windows. DSP and optional neural branches produce standardized synthetic-speech scores. Calibration, fusion, and sequential risk logic produce a decision tier.

\begin{figure}[t]
\centering
\fbox{\parbox{0.94\columnwidth}{\centering
PCM16/VoIP audio $\rightarrow$ validation and resampling $\rightarrow$ ephemeral buffer $\rightarrow$ VAD $\rightarrow$ 3-s windows\\[2pt]
$\swarrow$ DSP Fourier features \hspace{12mm} $\searrow$ optional neural detector\\[2pt]
calibrated score fusion $\rightarrow$ discounted Bayesian accumulation $\rightarrow$ low / medium / high event}}
\caption{High-level architecture of the proposed voice-cloning detection service.}
\label{fig:architecture}
\end{figure}

\section{DSP Fast Path}
\subsection{Fourier and Spectral Analysis}
For a window $x[n]$, the short-time Fourier transform is computed using a Hann window $w[n]$:
\begin{equation}
X(k,m)=\sum_{n=0}^{N-1}x[n+mH]w[n]e^{-j2\pi kn/N},
\end{equation}
where $N=512$ and $H=160$ for the reference 16-kHz configuration. The implementation derives spectral centroid, bandwidth, roll-off, flatness, flux, entropy, high-frequency energy ratio, zero-crossing rate, Mel-filterbank energies, and MFCC summaries. The Nyquist frequency is respected; high-frequency evidence is not treated as definitive when the source signal is narrowband.

Phase features include phase consistency and discontinuity statistics. More expensive group-delay and bicoherence features are configurable for offline experiments. Pitch is estimated using an autocorrelation proxy, and pause and voiced-activity ratios are acoustic proxies rather than verified respiration measurements.

\subsection{Voice Activity Detection and Windowing}
Short-time RMS energy is converted to decibels and compared with a configurable threshold. A window is processed only when its speech activity ratio exceeds the minimum. The ring buffer is bounded and expires stale sessions, preventing unbounded memory growth and preventing raw waveforms from entering operational stores.

\section{Neural Detector Architectures}
\subsection{CNN Baseline}
The baseline applies on-the-fly GPU STFT, logarithmic magnitude compression, and per-utterance normalization. The normalized tensor is passed to either a light convolutional network with Max-Feature-Map activations or a ResNet-18-style residual stack. Frequency is not aggressively pooled in the first block. Narrow frequency masking is applied during training as a configurable SpecAugment operation. The model returns logits $z\in\mathbb{R}^{B\times2}$ and an embedding $h\in\mathbb{R}^{B\times D}$.

\subsection{XLS-R and Graph-Attention Backend}
The main model uses the Hugging Face XLS-R checkpoint \texttt{facebook/wav2vec2-xls-r-300m}. Let $H_l$ denote the hidden state from layer $l$. Learned sensitive-layer selection produces:
\begin{equation}
H=\sum_{l=0}^{L-1}\operatorname{softmax}(a)_lH_l.
\end{equation}
The selected temporal sequence is processed by an AASIST-style graph-attention head that aggregates local temporal neighborhoods and performs attentive pooling. The reference configuration uses approximately $10^{-5}$ for the encoder and $2\times10^{-4}$ for the head. Feature-extractor freezing and layer truncation are configurable.

\section{Training and Calibration}
The training input is a CSV containing \texttt{location} and \texttt{label}; capitalization is accepted. Audio is loaded as mono, resampled to 16 kHz, and padded or truncated to three seconds. The training implementation supports weighted cross-entropy and an optional OC-Softmax objective. AdamW, warmup, cosine decay, gradient accumulation, gradient clipping, bf16/fp16 mixed precision, checkpoint resume, and validation-EER early stopping are configured through YAML.

Because classifier scores are not automatically calibrated probabilities, temperature scaling is fitted on a held-out validation set. Expected calibration error is reported using fixed probability bins. Thresholds are selected using bona fide validation scores and a target false-alarm rate, defaulting to one percent for the warning threshold. Temperature and thresholds are stored alongside the checkpoint and pinned at deployment.

\section{Sequential Risk Engine}
For a calibrated synthetic-speech probability $p_t$, the log-odds evidence is:
\begin{equation}
e_t=\log\left(\frac{p_t+\epsilon}{1-p_t+\epsilon}\right).
\end{equation}
The sequential state is updated as:
\begin{equation}
\ell_t=\ell_{t-1}+\lambda e_t,\qquad P_t=\sigma(\ell_t),
\end{equation}
where $\lambda\in(0,1]$ discounts evidence from overlapping and dependent windows. The discount, prior, and tier thresholds must be selected using validation data; independent-window multiplication is not assumed.

\section{Streaming Service and Privacy}
The FastAPI service exposes REST and WebSocket paths. The binary \texttt{/score} endpoint accepts exactly three seconds of PCM16 audio and returns a calibrated score, decision tier, model version, and latency. The WebSocket path accumulates speech windows and submits them to a bounded asynchronous batching queue. Model execution is moved off the event loop so multiple authorized calls can be served concurrently.

The service does not log raw waveforms or store audio in Redis, MongoDB, MLflow, or application logs. Metadata-only alerts may be retained with TTL policies. Memory-only processing does not guarantee physical erasure in the presence of swap, crash dumps, tracing, or host-level observability, which must be controlled operationally.

\section{Experimental Methodology}
\subsection{Dataset and Splits}
The dataset is represented by location-label records and may include speaker, source recording, generator, codec, accent, language, noise, SNR, and environment metadata. Primary experiments use speaker-disjoint and source-recording-disjoint partitions. Zero-day evaluation uses generator-disjoint and leave-one-generator-out tests. Cross-dataset, cross-accent, cross-codec, and noise robustness experiments are reported separately.

\subsection{Metrics}
We report equal error rate (EER), minimum detection cost function (minDCF), ROC-AUC, precision, recall, F1, false-alarm rate at 90\%, 95\%, and 99\% spoof recall, Brier score, expected calibration error, detection delay, false alarms per hour, throughput, memory, and p50/p95/p99 processing latency. Confidence intervals are computed by grouped bootstrap at the speaker, source, or call level. Test data are not used for threshold selection.

\subsection{Ablations}
Planned ablations compare DSP-only, neural-only, weighted fusion, calibrated fusion, and sequential fusion. DSP ablations include spectral, phase, pitch, Fourier, and pause-proxy subsets. Additional experiments vary window length, hop size, calibration, and evidence-discount factors.

\section{Results}
\textbf{This section is intentionally a placeholder until the authors run the real train, validation, test, and unseen-generator experiments.} Values from synthetic fixtures or debugging runs must not be reported as research evidence.

\begin{table}[t]
\centering
\caption{Primary held-out test results. Replace placeholders only with reproducible experiment outputs.}
\label{tab:primary-results}
\begin{tabular}{lcccccc}
\toprule
System & EER & minDCF & ROC-AUC & Brier & ECE & FA@99\\
\midrule
DSP-only & TBD & TBD & TBD & TBD & TBD & TBD\\
CNN baseline & TBD & TBD & TBD & TBD & TBD & TBD\\
XLS-R/AASIST & TBD & TBD & TBD & TBD & TBD & TBD\\
Hybrid sequential & TBD & TBD & TBD & TBD & TBD & TBD\\
\bottomrule
\end{tabular}
\end{table}

\section{Deployment Considerations}
The reference deployment uses a containerized FastAPI service on GPU-capable ECS EC2 instances behind an Application Load Balancer. Terraform variables control region, subnets, source security groups, instance type, artifact bucket, model checkpoint URI, and optional ACM certificate. A deployment circuit breaker enables rollback of unhealthy revisions. A managed SageMaker real-time endpoint is an alternative: ECS provides more direct control of long-lived WebSocket behavior and GPU capacity, while SageMaker reduces infrastructure management at the cost of endpoint constraints and managed-service coupling.

\section{Limitations and Ethical Considerations}
Acoustic artifacts vary by codec, device, language, accent, and recording environment. A detector may fail on unseen generators or produce false alarms for legitimate signal transformations. Pause and pitch features do not constitute medical or physiological measurements. High-risk scores should support human or policy-controlled verification rather than automatically denying service. The system requires privacy review, authorized deployment, and monitoring for demographic and channel-specific error differences.

\section{Conclusion}
This paper presents a reproducible architecture for detecting voice-cloning impersonation attacks in live VoIP calls. The design combines DSP evidence, trainable CNN and XLS-R/AASIST-compatible detectors, calibration, dependence-aware sequential risk, privacy-preserving streaming, and GPU deployment interfaces. The central empirical questions---generalization to unseen generators, calibration under channel shift, false alarms in realistic calls, and end-to-end latency---remain to be answered with the authors' real dataset and controlled experiments.

\section*{Acknowledgment}
The authors may add funding, laboratory, infrastructure, or dataset acknowledgments here.

\bibliographystyle{IEEEtran}
\begin{thebibliography}{00}
\bibitem{wav2vec2} A. Baevski, H. Zhou, A. Mohamed, and M. Auli, ``wav2vec 2.0: A framework for self-supervised learning of speech representations,'' in \emph{Proc. NeurIPS}, 2020.
\bibitem{xlsr} A. Conneau \emph{et al.}, ``Unsupervised cross-lingual representation learning for speech recognition,'' in \emph{Proc. Interspeech}, 2021.
\bibitem{aasist} K. Jung \emph{et al.}, ``AASIST: Audio anti-spoofing using integrated spectro-temporal graph attention networks,'' in \emph{Proc. ICASSP}, 2022.
\bibitem{asvspoof} ASVspoof Challenge, ``Automatic speaker verification spoofing and countermeasures,'' challenge documentation and evaluation protocols.
\bibitem{temperature} C. Guo, G. Pleiss, Y. Sun, and K. Q. Weinberger, ``On calibration of modern neural networks,'' in \emph{Proc. ICML}, 2017.
\end{thebibliography}

\end{document}
